\section{Architecture}


\subsection{Chain Placement}

Z-Chain operates as a specialized chain within the Zoo Network multi-chain architecture:

\begin{itemize}
\item \textbf{C-Chain} (Contract Chain): Application deployment. Consumes Z-Chain proofs via precompiles.
\item \textbf{Z-Chain} (ZK Substrate): Proof registration, verification, and receipt generation. This paper's focus.
\item \textbf{T-Chain} (Threshold Chain): FHE decryption coordination for private ZK-FHE hybrid protocols.
\item \textbf{K-Chain} (Key Chain): Post-quantum key management (ML-KEM, ML-DSA) for proof signing.
\end{itemize}

\subsection{Two-Lane Design}

\subsubsection{Production Lane}

Three proof systems are permanently supported with stable verifier contracts:

\begin{table}[h]
\centering
\begin{tabular}{llll}
\toprule
\textbf{System} & \textbf{Setup} & \textbf{PQ-Safe} & \textbf{Verify Gas} \\
\midrule
Groth16 & Trusted & No\textsuperscript{$\dagger$} & 230K \\
PLONK & Universal & No\textsuperscript{$\dagger$} & 350K \\
STARK & Transparent & Yes & 500K \\
\bottomrule
\end{tabular}
\caption{Production proof systems. $\dagger$: Wrapped in PQZ for quantum resistance.}
\end{table}

\subsubsection{Research Lane}

Experimental proof systems register under versioned capability IDs ($\geq 100$). Each carries an explicit risk label (\texttt{EXPERIMENTAL}, \texttt{AUDITED}, \texttt{DEPRECATED}) and runs in sandboxed verifier contracts that cannot affect production lane state.

